% SPDX-License-Identifier: LPPL-1.3c
% Copyright (C) 2026 Christophe Jorssen
% Author and Current Maintainer: Christophe Jorssen <christophe.jorssen@gmail.com>
% LPPL maintenance status: maintained
% This file is part of luafemm. See LICENSE and LICENSES.md for its terms.
\documentclass[varwidth=155mm,border=3mm]{standalone}
\usepackage{tikz,pgfplots}
\usetikzlibrary{femm.circuits.ee.IEC}
\usepgfplotslibrary{femm}
\pgfplotsset{compat=1.18}
\begin{document}

\begin{tikzpicture}[
  femm/problem={name=network-demo,frequency=200,
    xmin=-10,xmax=10,ymin=-8,ymax=8,mesh size=1,depth=20},
  femm/circuit={name=coil,drive=network},
  femm/network={name=supply},
  femm/element={name=source,kind=voltage source,p=a,n=0,value=1},
  femm/element={name=load,kind=resistor,p=a,n=b,value=2},
  femm/element={name=winding,kind=winding,p=b,n=0,circuit=coil},
  scale=2]
  \filldraw[fill=orange!25,
    femm/region={material=copper,circuit=coil,turns=20}]
    (-6,-4) rectangle (-4,4);
  \filldraw[fill=orange!25,
    femm/region={material=copper,circuit=coil,turns=-20}]
    (4,-4) rectangle (6,4);
  \femmmeasure[name=coil-terminal,quantity=circuit,circuit=coil]
  \femmmeasure[name=coil-element,quantity=network,element=winding]
  \pic[femm/lines=10] {femm field};
\end{tikzpicture}

\par

\begin{tikzpicture}[circuit ee IEC]
  \draw (0,0) to[voltage source={info={$1\,\mathrm V$}}] (0,2)
    to[resistor={info={$2\,\Omega$}}] (4,2)
    to[inductor={femm/ee={model=network-demo,
      element=winding,result=coil-element,values=solution,
      show=current,representation=rms}}] (4,0) -- (0,0);
\end{tikzpicture}

\par

\begin{tikzpicture}
  \pic[femm/ee view={model=network-demo}] {femm ee circuit};
\end{tikzpicture}

\end{document}
